\documentclass[10pt,a4paper]{amsart}
\usepackage[margin=19mm]{geometry}
\usepackage{amsmath,amssymb,mathtools}
\usepackage[scr=rsfs,cal=euler]{mathalfa}
\usepackage{xfrac}
\usepackage[unicode,hidelinks,bookmarks=false]{hyperref}
\newcommand{\ie}{i.e.\ }
\newcommand{\cf}{cf.\ }
\newcommand{\resp}{respectively\ }
\newcommand{\Subsection}{\subsection}
\providecommand{\nobreakdash}{\nobreak\hbox{-}}
\setlength{\emergencystretch}{2em}
\multlinegap0pt
\usepackage[normalem]{ulem}
\usepackage{enumerate}
\usepackage{mathtools}
\usepackage{amssymb}
\newtheorem{lem}{Lemma}
\usepackage{cleveref}
\crefname{lem}{Lemma}{Lemmas}
\renewcommand{\deg}{\mathsf{deg}}
\renewcommand{\geq}{\geqslant}
\renewcommand{\leq}{\leqslant}
\renewcommand{\max}{\mathsf{max}}
\newcommand\mc\mathcal
\newcommand\nat{\mathbb N}
\title{Arithmetic genus inequalities with an application to sums of squares}
\author{David Grimm, Gonzalo Manzano-Flores}
\date{Source: arXiv:2501.03369v2 (CC BY 4.0)}
\begin{document}
\maketitle

\noindent\emph{Source and licence.} Arithmetic genus inequalities with an application to sums of squares. Version arXiv:2501.03369v2 (CC BY 4.0), distributed by arXiv under the Creative Commons Attribution 4.0 International licence, http://creativecommons.org/licenses/by/4.0/, as recorded in the arXiv licence metadata for this exact version and shown on the abstract page https://arxiv.org/abs/2501.03369v2. Full licence notice: \texttt{licenses/arxiv-2501.03369-CC-BY-4.0.txt}.\par\smallskip

\noindent\textit{Editorial scope.} Counted unit: the article's lem L:G-fixpoint-pivot-count (source0.tex lines 262--265). The article's complete original proof of it is delivered with it (source0.tex lines 267--297, one \texttt{proof} environment of 31 lines). No mathematics is written, repaired or re-derived: the statement and the proof are the article's own lines, in the article's own notation, and no retained line is substituted or re-flowed.  Retained uncounted as the source-local context the proof needs: lines 237--239.  Nothing was omitted from any retained span, so the package carries no omission record.  No retained line contains a citation, so the wrapper carries no bibliography and no bibliography entry.  No retained line contains a figure environment, an image-inclusion command or drawing code, so no image is delivered and the package declares no asset dependency.  Editorial formatting only, in addition: an AMS \texttt{amsart} preamble that restates the article's own declarations and macros used by the retained lines, namely the environments \texttt{lem} on separate counters, together with the standard packages those lines need.  The retained environments print in this document as Lemma 1 in order of appearance, whereas the article prints the same items with the numbering of its own full document.  The complete native source of the article as distributed in the arXiv e-print for this exact version is bundled unchanged as \texttt{sources/171236/source0.tex}.
\begin{lem}\label{restricting rigidities}
   Let $G$ be a group and $\mathcal{D}=(\mathcal{V},\mathcal{E})$ a finite simple $G$-graph. Let $\mathcal{D}'=(\mathcal{V}',\mathcal{E}')$ be a connected $G$-stable subgraph of $\mathcal{D}$. Let $\mathcal{R}=(\mathcal{V}_\mathcal{R},\mathcal{E}_\mathcal{R})$ be a rigidity of the $G$-graph $\mathcal{D}$. Then every connected component of $\mathcal{R}':=(\mathcal{V}_{\mathcal{R}} \cap \mathcal{V}' ,  \mathcal{E}_\mathcal{R}\cap \mathcal{E}')$ is a rigidity of the $G$-graph $\mathcal{D}'$.
\end{lem}

\begin{lem}\label{L:G-fixpoint-pivot-count}
Let $\mc{D}^G_G\neq \emptyset$.
Then $|\mc{D}_G | \leq \beta(\mc{D}) + 1$.  
\end{lem}

\begin{proof}
Let  $\mc{R}_0 \in \mc{D}^G_G$. We will show that  $|\mc{D}_G \setminus \{\mc{R}_0\} | \leq \beta(\mc{D})$.
For $v\in\mc{V}$ let $d(v)$ denote the distance between $v$ and $\mc{R}_0$, that is,
the smallest $r\in\nat$ such that there exist $v_0,v_1,\ldots,v_r\in\mc{V}$ with $\{v_{i-1},v_i\}\in\mc{E}$ for $1\leq i\leq r$ where $v_r=v$ and $v_0 \in \mc{R}_0$.
Note that $d(gv)=d(v)$ for all $v\in\mc{V}$ and $g\in G$.

The proof of the statement is  by induction 
 on $|\mc{E}|$.
If $|\mc{E}|=0$, then $\mc{V}= \mc{R}_0=\{v_0\}$ and the statement holds trivially.
Assume now that $|\mc{E}|\geq 1$. We set $m=\max\{d(v)\mid v\in\mc{V}\}$, $\mc{M}=\{v\in\mc{V}\mid d(v)=m\}$ and $\mc{E}^\ast=\mc{E}\cap\{\{w,w'\}\mid w,w'\in\mc{M}\}$.

We first consider the case where $\mc{E}^\ast\neq \emptyset$.
We set $\mc{E}'=\mc{E}\setminus\mc{E}^\ast$ and consider the graph $\mc{D}'=(\mc{V},\mc{E}')$.
Note that $\mc{D}'$ is connected, $\beta(\mc{D}') <\beta(\mc{D})$, and the $G$-action on $\mc{D}$ restricts to a $G$-action on $\mc{D}'$.
Furthermore, every rigidity of $\mc{D}$, after removing edges from $\mc{E}^\ast$, either remains connected and hence remains a single rigidity of $\mc{D}'$, or it decomposes in several connected components, each of which gives a rigidity on $\mc{D}'$, by \Cref{restricting rigidities}.
Since $|\mc{E}'|<|\mc{E}|$, we therefore obtain by the induction hypothesis that 
$|\mc{D}_G\setminus\{R_0\}|\leq |\mc{D}'_G\setminus\{R_0\}|\leq \beta(\mc{D}')< \beta(\mc{D})$.

We now consider the case where $\mc{E}^\ast=\emptyset$. 
Set $\mc{V}'=\mc{V}\setminus \mc{M}$. We consider the graph $\mc{D}'=(\mc{V}',\mc{E}')$ where $\mc{E}'=\mc{E}\cap\{\{v,v'\}\mid v,v'\in\mc{V}'\}$ (i.e.~the full subgraph of $\mc{D}$ spanned by $\mc{V}'$).
Note that $\mc{D}'$ is connected,  $\beta(\mc{D}')\leq \beta(\mc{D})$, the $G$-action on $\mc{D}$ restricts to a $G$-action on $\mc{D}'$, and every rigidity of $\mc{D}$ that contains a vertex in $\mc{V}'$ restricts to at least one rigidity of $\mc{D}'$. 
Since $\mc{M}\neq \emptyset$ we have that $|\mc{E}'|<|\mc{E}|$.
Hence, the induction hypothesis yields that  $|\mc{D}'_G\setminus\{R_0\}|\leq \beta(\mc{D}')$. Let $\mc{M}_G$ denote the set of (necessarily singular) rigidities of $\mc{D}_G$ that contain no vertex in $\mc{V}'$.  Then 
$\vert \mc{D}_G\setminus\{\mc{R}_0\} \vert \leq \vert \mc{D}'_G\setminus\{\mc{R}_0\} \vert + \vert \mc{M}_G \vert \leq \beta(\mc{D}'_G) + \vert \mc{M}_G  \vert$.
Since $\mc{E}^\ast=\emptyset$, we have $|\mc{E}\setminus\mc{E}'| = \sum_{w \in \mc{M}} \deg(w)$ and therefore
$\beta(\mc{D})-\beta(\mc{D}')=\sum_{w\in\mc{M}}(\deg(w)-1)$.
Since every vertex in $\mc{M}_G$ has degree at least $2$, we obtain that 
 $\beta(\mc{D})-\beta(\mc{D}')\geq |\mc{M}_G|$.
Hence we conclude that 
 $|\mc{D}_G\setminus\{R_0\}|\leq |\mc{D}'_G \setminus\{R_0\}|+|\mc{M}_G|\leq \beta(\mc{D})$.
\end{proof}

\end{document}
